\documentclass[a4paper,twoside,12pt]{article}
\usepackage[top=1.5cm,left=1.5cm,right=1.5cm,bottom=2cm]{geometry}

\usepackage[utf8]{inputenc}
\usepackage[T1]{fontenc}
\usepackage{lmodern}
\usepackage[brazilian]{babel}

\usepackage{graphicx}
\usepackage{subfig}
\usepackage{psfrag}
\usepackage{cite}
\usepackage[cmex10]{amsmath}
\usepackage{amssymb}
\usepackage{upgreek}
\usepackage{microtype}
\usepackage{titling}
\usepackage{courier}
\usepackage{url}
\usepackage{hyperref}
\hypersetup{
	colorlinks=true,
	linkcolor=blue,
	urlcolor=blue
}


% Use the 'listings' package to add code and the 'color' package to generate the highlights
\usepackage{listings}
\usepackage{color}
\definecolor{grey}{rgb}{0.6,0.6,0.6}
\definecolor{codegreen}{rgb}{0,0.6,0}
\definecolor{codegray}{rgb}{0.5,0.5,0.5}
\definecolor{codepurple}{rgb}{0.58,0,0.82}
\definecolor{backcolour}{rgb}{0.95,0.95,0.92}
\lstset{language=Python,				% set it to Python language highlighting
	backgroundcolor=\color{backcolour},
	commentstyle=\color{codegreen},
	keywordstyle=\color{magenta},
	numberstyle=\tiny\color{codegray},
	stringstyle=\color{codepurple},
	basicstyle=\ttfamily\small,		% set size of fonts
	breaklines=true,				% set automatic line breaking
	linewidth=\textwidth,			% set size of code box
	showstringspaces=false}			% show spaces as underscores only inside strings



\title{\vspace{-2cm} Processamento Digital de Sinais e Aplicações em Acústica\\ Tutorial 07 - Série de Fourier e a Transformada Discreta de Fourier}
\author{\url{https://github.com/fchirono/Aulas_PDS_Acustica}}

\begin{document}
	\date{}
	\maketitle
	
	\section*{Objetivos do Tutorial}
	Ao final desta sessão, você será capaz de:
	\begin{itemize}
		\item Identificar as condições necessárias para que a DFT de um sinal discreto resulte exatamente nos coeficientes da Série de Fourier,
		\item Compreender a relação entre os parâmetros de aquisição (frequência de amostragem, duração do sinal, frequência fundamental) e a qualidade dos coeficientes obtidos via DFT,
		\item Sintetizar sinais periódicos a partir de coeficientes da Série de Fourier e verificar sua reconstrução no domínio do tempo,
		\item Calcular e interpretar a DFT de sinais periódicos, comparando os resultados numéricos com os coeficientes teóricos esperados.
	\end{itemize}
	
	\begingroup
	\let\clearpage\relax
	\tableofcontents
	\endgroup
	
	
	% -----------------------------------------------------------------------
	\section{Introdução}
	\label{sec:intro}
	
	A Série de Fourier é uma representação matemática fundamental de sinais periódicos, decompondo-os em uma soma de senoides de frequências múltiplas de uma frequência fundamental $f_0$. Na prática, o processamento digital de sinais lida com versões discretizadas desses sinais, e a ferramenta numérica adequada para calcular os coeficientes da Série de Fourier de um sinal discreto é a \emph{Transformada Discreta de Fourier} (DFT), implementada eficientemente pelo algoritmo da \emph{Transformada Rápida de Fourier} (FFT).
	
	Este tutorial explora a relação entre a Série de Fourier e a DFT, apresentando as equações de síntese e de análise da série, e destacando as condições necessárias para que a DFT produza exatamente os coeficientes da Série de Fourier.
	
	\subsection{A Série de Fourier -- Equações de Síntese}
	
	Um sinal periódico $x(t)$ com período $T_0 = 1/f_0$ pode ser representado pela Série de Fourier. As \emph{equações de síntese} reconstroem o sinal a partir dos seus coeficientes. Na forma real (seno e cosseno):
	
	\begin{equation}
		x(t) = \sum_{k=0}^{\infty} \left[ A_k \cos(2\pi k f_0 t) + B_n \sin(2\pi k f_0 t) \right],
		\label{eq:sintese_real}
	\end{equation}
	
	\noindent onde $A_k$ e $B_k$ são os coeficientes reais das componentes de cosseno e seno do $k$-ésimo harmônico, respectivamente. Na forma exponencial complexa:
	
	\begin{equation}
		x(t) = \sum_{k=-\infty}^{+\infty} C_k \, e^{\,j \, 2\pi k f_0 t},
		\label{eq:sintese_complexa}
	\end{equation}
	
	\noindent onde $C_k$ são os coeficientes complexos da série exponencial. A relação entre os dois conjuntos de coeficientes é:
	
	\begin{equation}
		C_k = \frac{A_k}{2} - j\frac{B_k}{2}, \qquad C_{-k} = \frac{A_k}{2} + j\frac{B_k}{2}, \qquad k \geq 1,
		\label{eq:relacao_coefs}
	\end{equation}
	
	\noindent e $C_0 = A_0$ para o termo DC ($k = 0$).
	
	\subsection{A Série de Fourier -- Equações de Análise}
	
	As \emph{equações de análise} permitem calcular os coeficientes diretamente a partir do sinal $x(t)$, explorando a ortogonalidade das funções seno e cosseno no intervalo $[0, T_0]$. Os coeficientes reais são obtidos por:
	
	\begin{align}
		A_0 &= \frac{1}{T_0} \int_0^{T_0} x(t) \, dt,
		\label{eq:analise_A0} \\[6pt]
		A_k &= \frac{2}{T_0} \int_0^{T_0} x(t) \cos(2\pi k f_0 t) \, dt, \quad k \geq 1,
		\label{eq:analise_Ak} \\[6pt]
		B_k &= \frac{2}{T_0} \int_0^{T_0} x(t) \sin(2\pi k f_0 t) \, dt, \quad k \geq 1.
		\label{eq:analise_Bk}
	\end{align}
	
	Note que $A_0$ corresponde ao valor médio (componente DC) do sinal, e o fator $2/T_0$ (em lugar de $1/T_0$) nos coeficientes $A_k$ e $B_k$ com $k \geq 1$ resulta da normalização pela metade da energia da senoide no período. Na forma exponencial, o coeficiente complexo do $k$-ésimo harmônico é dado por:
	
	\begin{equation}
		C_k = \frac{1}{T_0} \int_0^{T_0} x(t) \, e^{-j2\pi k f_0 t} \, dt, \quad k \in \mathbb{Z}.
		\label{eq:analise_Ck}
	\end{equation}
	
	\noindent Esta expressão é válida para todo $k$ inteiro (positivo, negativo e zero), e unifica as equações \eqref{eq:analise_A0}--\eqref{eq:analise_Bk} em uma única fórmula. A integral pode ser calculada sobre qualquer intervalo de duração $T_0$; o intervalo $[0, T_0]$ é utilizado por conveniência.
	
	\subsection{A DFT e os Coeficientes da Série de Fourier}
	
	A DFT de um vetor de $N_t$ amostras $x[n]$, com $n = 0, 1, \ldots, N_t-1$, é definida como:
	
	\begin{equation}
		X[k] = \sum_{n=0}^{N_t-1} x[n] \, e^{-j2\pi kn/N_t}, \quad k = 0, 1, \ldots, N_t-1.
		\label{eq:dft}
	\end{equation}
	
	As frequências discretas $k$ correspondem às frequências analógicas (em Hertz):
	
	\begin{equation}
		f_k = k \, \frac{f_s}{N_t} = k \, \Delta f,
		\label{eq:freq_discreta}
	\end{equation}
	
	\noindent onde $f_s$ é a frequência de amostragem e $\Delta f = f_s / N_t$ é o intervalo entre duas frequências discretas adjacentes. 
	
	% -----------------------------------------------------------------------
	\section{Tarefas}
	\label{sec:tarefas}
	
	Inicie o seu script Python com as seguintes declarações iniciais:
	
	\begin{lstlisting}[frame=single, caption={Configurações iniciais comuns aos exercícios.}]
import numpy as np
import matplotlib.pyplot as plt
plt.close("all")

# Frequencia de amostragem [Hz]
fs  = 10000
dt  = 1/fs          # intervalo de amostragem [s]

# Frequencia fundamental [Hz]
f0  = 100.
T   = 1/f0          # duracao de um periodo [s]
Nt  = int(T*fs)     # No. de amostras (= 1 periodo exato)
t   = np.linspace(0, T-dt, Nt)  # vetor de tempo

# Numero de coeficientes da Serie de Fourier
K = 25
	\end{lstlisting}
	
	
	% -----------------------------------------------------------------------
	\subsection{Correspondência entre DFT e Série de Fourier}
	\label{sec:tarefa_condicoes}
	
		
	\begin{enumerate}
		\item Descreva analiticamente as condições necessárias para que a DFT de um sinal discreto produza exatamente os coeficientes da Série de Fourier. 
	\end{enumerate}
	 
	Para os exercícios seguintes, utilize como sinal de referência a onda dente-de-serra, cujos coeficientes de seno são dados por:
	
	\begin{equation}
		B_k = -\frac{2(-1)^{k+1}}{\pi n}, \quad k = 1, 2, 3, \ldots
		\label{eq:coef_serra}
	\end{equation}
	
	\begin{enumerate}	
		\setcounter{enumi}{1}
		\item Usando as configurações iniciais acima, sintetize o sinal $x[n]$ no domínio do tempo discreto somando $K = 25$ termos da Série de senos de Fourier (equação de síntese \eqref{eq:sintese_real} com $A_k = 0$):
		\begin{equation*}
			x[n] = \sum_{k=1}^{K} B_k \sin(2\pi k f_0 n).
		\end{equation*}
		Plote o sinal sintetizado no domínio do tempo. O sinal obtido se assemelha a uma onda dente-de-serra? O que você observa ao usar mais termos na série?
		
		\item Calcule a DFT do sinal usando \texttt{numpy.fft.fft} e normalize-a pelo número de amostras $N_t$. Plote o espectro de magnitude e fase em função da frequência em Hz. Compare visualmente os coeficientes obtidos pela DFT com os coeficientes teóricos $C_n$ calculados a partir de $B_n$ via equação \eqref{eq:relacao_coefs}. Note que os coeficientes $C_{-n}$ aparecem nos índices superiores do vetor da DFT (isto é, as frequências negativas são espelhadas para as frequências acima de $f_s/2$).
		
		\item Experimente com uma duração \textbf{não inteira}: use $T_{ni} = 1{,}5 / f_0$ (1,5 períodos). Sintetize o sinal com essa duração e calcule sua DFT. O que acontece com os coeficientes em comparação com os casos anteriores? Explique o resultado observado com base nas condições descritas na questão 1.
		
	\end{enumerate}
	
	
	% -----------------------------------------------------------------------
	\subsection{Síntese e Análise de Diferentes Formas de Onda}
	
	Nesta tarefa, você irá trabalhar com quatro formas de onda periódicas clássicas: onda dente-de-serra, onda quadrada, e duas formas de onda triangular. As funções para calcular os coeficientes analíticos de cada onda são fornecidas abaixo:
	
	\begin{lstlisting}[frame=single, caption={Funções para calcular coeficientes analíticos de diferentes formas de onda.}]
def coef_dente_sen(m):
	# Onda dente-de-serra (coeficientes de seno)
	return -2*((-1)**(m+1))/(np.pi*m)

def coef_quad_sen(m):
	# Onda quadrada (coeficientes de seno)
	return 2*(1 - np.cos(np.pi*m))/(m*np.pi)

def coef_triang_cos(m):
	# Onda triangular, x em [-1, 1] (coeficientes de cosseno)
	return 8*np.sin(m*np.pi/2) / (np.pi*m)**2

def coef_triang_sen(m):
	# Onda triangular, x em [0, pi] (coeficientes de seno, com DC)
	return 2*(np.cos(np.pi*m) - 1)/(np.pi*m**2)
	\end{lstlisting}
	
	\noindent Use as configurações iniciais da Seção~\ref{sec:tarefas} com $N = 25$ termos. Repita os itens abaixo para cada uma das quatro formas de onda.
	
	\begin{enumerate}
		
		\item Sintetize o sinal no domínio do tempo usando a equação de síntese \eqref{eq:sintese_real} e plote-o. O formato visual do sinal converge para a forma de onda esperada com $N = 25$ termos?
		
		\item Calcule os coeficientes teóricos $C_n$ e $C_{-n}$ a partir de $A_n$ e $B_n$, usando as relações da equação \eqref{eq:relacao_coefs}. Verifique também se os coeficientes analíticos fornecidos acima são consistentes com as equações de análise \eqref{eq:analise_Ak}--\eqref{eq:analise_Bk}: para a onda quadrada, por exemplo, calcule numericamente a integral $B_n = \frac{2}{T_0}\int_0^{T_0} f(t)\sin(2\pi n f_0 t)\,dt$ usando \texttt{numpy.trapz} e compare com \texttt{coef\_quad\_sen(n)}.
		
		\item Calcule a DFT do sinal sintetizado e compare, em um mesmo gráfico de magnitude e fase, os coeficientes obtidos numericamente pela DFT com os coeficientes teóricos $C_n$. A correspondência é exata?
		
		\item Para a onda quadrada, observe que apenas os harmônicos ímpares estão presentes na Série de Fourier. Verifique se a DFT reproduz corretamente este comportamento, e confirme que os coeficientes dos harmônicos pares são numericamente próximos de zero.
		
	\end{enumerate}
	
	
	
	% -----------------------------------------------------------------------
	\section*{Referências}
	
	\begin{enumerate}
		\item K. Shin e J. Hammond, \textit{Fundamentals of Signal Processing for Sound and Vibration Engineers}. John Wiley and Sons, 2008.
		\item A. V. Oppenheim e R. W. Schafer, \textit{Discrete-Time Signal Processing} (3a Ed), Prentice Hall, 2010.
		\item J. G. Proakis e D. G. Manolakis, \textit{Digital Signal Processing: Principles, Algorithms and Applications} (4a Ed), Pearson, 2007.
	\end{enumerate}
	
\end{document}